\section{Averaging and completion of the proof}
\label{sec:conclusion}

We first derive the lower bound that contradicts the chord estimate of
Section~\ref{sec:chords}. Throughout this argument, \(n>2\), and we retain
the normalized extremizing cycle \(A\), its mass measure \(\mu\), and the
nonnegative minimizer \(v\) of Proposition~\ref{prop:minimizer}. Thus
\[
 \beta=\frac1{n-2},\qquad p=\frac{2n}{n-2}=2+4\beta,
 \qquad W=\int_Y v^p\dd\mu>0,
 \qquad \dd\nu=\frac{v^p}{W}\dd\mu.
\]
We use the same finite, nonnegative Borel representative of \(v\) as in
Section~\ref{sec:chords}.

\begin{lemma}[Opposite averaged chord bound]
\label{lem:opposite-average}
Define
\begin{equation}
 \label{finish:weighted-moment}
 B_*:=\int_Y v^{p+2\beta}\dd\mu,
 \qquad
 I(y):=\int_Y d(y,x)^2v(x)^{p+2\beta}\dd\mu(x).
\end{equation}
Then \(0<B_*<\infty\), the function \(I\) is finite and continuous on
\(Y\), and
\begin{equation}
 \label{finish:all-centers}
 W^2\le B_*I(y)\qquad\text{for every }y\in Y.
\end{equation}
In particular, for the chord moments
\[
 L_y=\int_Y
       \bigl[v(y)^\beta d(y,x)v(x)^\beta\bigr]^2\dd\nu(x),
\]
one has
\begin{equation}
 \label{finish:opposite-average}
 \int_Y L_y\dd\nu(y)\ge2.
\end{equation}
\end{lemma}

\begin{proof}
Lemma~\ref{lem:powers} gives all finite moments of \(v\), so \(B_*\) is
finite. It is positive because \(W>0\). Write \(D_Y=\diam Y\). Then
\(0\le I(y)\le D_Y^2B_*\), and the triangle inequality gives
\[
 |I(y)-I(z)|\le2D_YB_*d(y,z)\qquad(y,z\in Y).
\]
This proves the stated continuity.

Fix any \(y\in Y\) and put \(r=d(y,\cdot)\). By
Lemma~\ref{lem:powers}, the nonnegative function \(g=v^p\) belongs to
\(\Hcal\), and \(Dg=pv^{p-1}Dv\). The extension of the radial inequality
\eqref{eq:radial} furnished by Proposition~\ref{prop:sobolev-space}
therefore applies to this \(g\). Dividing by \(n\) and using
\(p/n=2\beta\) gives
\begin{equation}
 \label{finish:radial-test}
 W\le\int_Y r\left(v^p\sqrt{1-|Dr|^2}
             -2\beta v^{p-1}\ip{Dr}{Dv}\right)\dd\mu.
\end{equation}
At almost every point of each current chart, take the direct sum of its
cotangent inner-product space with \(\R\). The vector
\[
 \left(\sqrt{1-|Dr|^2},Dr\right)
\]
has norm one. Moreover, the identities
\[
 \frac p2+\beta+1+\beta=p,
 \qquad
 \frac p2+\beta+\beta=p-1
\]
give the factorization
\begin{align*}
 &v^p\sqrt{1-|Dr|^2}-2\beta v^{p-1}\ip{Dr}{Dv}\\
 &\quad=
 v^{p/2+\beta}
 \left\langle
   \left(\sqrt{1-|Dr|^2},Dr\right),
   \left(v^{1+\beta},-2\beta v^\beta Dv\right)
 \right\rangle\\
 &\quad\le
 v^{p/2+\beta}
 \left(v^{2+2\beta}+4\beta^2v^{2\beta}|Dv|^2\right)^{1/2}.
\end{align*}
All powers in this formula are nonnegative powers of \(v\), and the
identity also holds on \(\{v=0\}\). Applying the integral
Cauchy--Schwarz inequality to \eqref{finish:radial-test} now yields
\begin{equation}
 \label{finish:cauchy}
 W^2\le I(y)
 \int_Y\left(v^{2+2\beta}
             +4\beta^2v^{2\beta}|Dv|^2\right)\dd\mu.
\end{equation}
Both factors are finite: the first was bounded above, and the second is
finite by the moment estimates and the Sobolev power
\(v^{1+\beta}\) in Lemma~\ref{lem:powers}.

To identify the last integral, use
\(\psi=v^{1+2\beta}\in\Hcal\) in the Euler equation
\eqref{eq:euler}. Its gradient is
\(D\psi=(1+2\beta)v^{2\beta}Dv\). Thus
\[
 4\beta(1+2\beta)\int_Yv^{2\beta}|Dv|^2\dd\mu
 +n\int_Yv^{2+2\beta}\dd\mu
 =n\int_Yv^{p+2\beta}\dd\mu.
\]
Since \(1+2\beta=n\beta\), division by \(n\) proves
\begin{equation}
 \label{finish:euler-moment}
 \int_Y\left(v^{2+2\beta}
             +4\beta^2v^{2\beta}|Dv|^2\right)\dd\mu=B_*.
\end{equation}
Equations~\eqref{finish:cauchy} and \eqref{finish:euler-moment}
prove \eqref{finish:all-centers}. The argument fixed an arbitrary center
before taking chart derivatives. Hence its conclusion holds for every
center, although the null set on which \(Dr\) is undefined may depend
on that center.

We have proved the all-center estimate \(W^2\le B_*I(y)\).
It remains to combine it with the CAT(0) variance inequality
\cite[Proposition~4.4]{Sturm2003}, which doubles this lower contribution
when we average over centers. We recall its proof for the finite measure
$v^{p+2\beta}\mu$.
Compactness of \(Y\) and continuity of \(I\) provide a minimizer
\(o\in Y\). This is a barycenter of the finite measure
\(v^{p+2\beta}\mu\). For \(y\in Y\), let \(o_t\) be the point at
fraction \(t\in(0,1)\) on the segment from \(o\) to \(y\). The
CAT(0) squared-distance inequality, integrated against
\(v(x)^{p+2\beta}\dd\mu(x)\), gives
\[
 I(o)\le I(o_t)
 \le(1-t)I(o)+tI(y)-t(1-t)B_*d(o,y)^2.
\]
Subtracting \((1-t)I(o)\), dividing by \(t\), and letting
\(t\downarrow0\) proves the variance inequality
\begin{equation}
 \label{finish:variance}
 I(y)\ge I(o)+B_*d(o,y)^2\qquad(y\in Y).
\end{equation}
This uses only the segment from \(o\) to \(y\).

The chord kernel is a nonnegative Borel function on \(Y\times Y\).
Its integral is finite, since its expression below is bounded above by
\(D_Y^2B_*^2/W^2\). Tonelli's theorem and
\eqref{finish:weighted-moment} give
\begin{equation}
 \label{finish:double-integral}
 \begin{aligned}
 \int_YL_y\dd\nu(y)
 &=\iint_{Y\times Y}
       \bigl[v(y)^\beta d(y,x)v(x)^\beta\bigr]^2
             \dd\nu(x)\dd\nu(y)\\
 &=\frac1{W^2}\int_Y I(y)v(y)^{p+2\beta}\dd\mu(y)\\
 &\ge\frac1{W^2}\left(
       B_*I(o)+B_*\int_Y d(o,y)^2v(y)^{p+2\beta}\dd\mu(y)
                  \right)\\
 &=\frac{2B_*I(o)}{W^2}\ge2.
 \end{aligned}
\end{equation}
Here the first inequality is \eqref{finish:variance}, and the last is
\eqref{finish:all-centers} at \(y=o\). The averaged expressions contain
only distances and values of the Borel representative of \(v\); the
proof requires no joint choice of the derivatives of distance
functions. This proves \eqref{finish:opposite-average}.
\end{proof}

\begin{proof}[Proof of Theorem~\ref{thm:main}]
We first prove the statement in every compact CAT(0) space by induction
on the cycle dimension \(n\ge2\). For such a space \(Y\), let
\(V(T)\) denote the infimum of the masses of integral fillings of an
integral \(n\)-cycle \(T\) in \(Y\). Theorem~\ref{thm:current-background}
ensures that this infimum is finite and attained. It therefore suffices
in each dimension to prove
\begin{equation}
 \label{finish:compact-bound}
 V(T)\le c_n\M(T)^{(n+1)/n}
 \qquad\text{for every integral }n\text{-cycle }T\text{ in }Y.
\end{equation}
The zero cycle has the zero current as a filling.

In dimension \(n=2\), a violation of
\eqref{finish:compact-bound} would, by Lemma~\ref{lem:extremizer},
produce a nonzero extremizing cycle that can be rescaled to satisfy
\eqref{eq:normalization}. Proposition~\ref{prop:base} excludes this
cycle: its radial inequality and Euclidean lower density force
\(\M(A)\ge s_2\), whereas the normalization requires
\(\M(A)<s_2\). Thus \eqref{finish:compact-bound} holds in
dimension two, and attainment gives an integral filling with that mass
bound.

Now let \(n>2\), and assume the asserted existence of sharp integral
fillings in dimension \(n-1\) in every compact CAT(0) space. Suppose
that \eqref{finish:compact-bound} fails for a cycle \(T\) in a compact
CAT(0) space \(Y\). Since \(T\ne0\), we may choose
\[
 c_n<c<\frac{V(T)}{\M(T)^{(n+1)/n}}.
\]
Lemma~\ref{lem:extremizer} produces a nonzero cycle \(A\) maximizing
\(V(B)-c\M(B)^{(n+1)/n}\) over integral \(n\)-cycles. After a constant
rescaling of the metric it satisfies \eqref{eq:extremizer} and
\eqref{eq:normalization}, in particular
\[
 m:=\M(A)=((n+1)c)^{-n}<s_n.
\]
The rescaled space remains compact and CAT(0), so the induction
hypothesis continues to apply there.

Proposition~\ref{prop:radial} gives \eqref{eq:radial} for this cycle.
The sharp filling statement in dimension \(n-1\) is exactly the
induction input to Proposition~\ref{prop:small-set}. Consequently
Section~\ref{sec:sobolev}, and in particular the almost sharp Sobolev
inequality of Proposition~\ref{prop:almost-sobolev}, applies.
Propositions~\ref{prop:sobolev-space} and
\ref{prop:compact-embedding} supply the completed calculus and
compactness used in Proposition~\ref{prop:minimizer}. We therefore
obtain a nonnegative minimizer \(v\) satisfying \eqref{eq:euler} and
\eqref{eq:minimizer-normalization}, with
\[
 0<W=\int_Yv^p\dd\mu<s_n.
\]
Lemmas~\ref{lem:powers} and \ref{lem:composites} justify the powers
and nonlinear tests used in Section~\ref{sec:chords} and in
Lemma~\ref{lem:opposite-average}.

For the probability measure \(\nu=v^p\mu/W\), the chord upper bound
\eqref{eq:chord-upper} holds for \(\nu\)-almost every center. Integrating
that bound and using Lemma~\ref{lem:opposite-average} gives
\[
 2\le\int_YL_y\dd\nu(y)
   \le2\left(\frac W{s_n}\right)^{2/n}<2,
\]
a contradiction. Hence no violating cycle \(T\) exists, and
\eqref{finish:compact-bound} holds in dimension \(n\). By attainment
in Theorem~\ref{thm:current-background}, every integral \(n\)-cycle
has an integral filling realizing \(V(T)\), with the required mass
bound. Its support is compact because it is a closed subset of \(Y\).
This completes the induction for every \(n\ge2\).

Lemma~\ref{lem:compact-reduction} now transfers this bound to every
compactly supported integral cycle \(T\in\I_n(X)\) in a proper
CAT(0) space. It supplies a compactly supported integral filling
\(S\) with
\[
 \partial S=T,\qquad
 \M(S)\le c_n\M(T)^{(n+1)/n}
 =\frac{\M(T)^{(n+1)/n}}
 {(n+1)^{(n+1)/n}\omega_{n+1}^{1/n}}.
\]
The zero cycle has the zero filling. This proves the theorem.
\end{proof}
